\section{Progress as of 24 August 2026}
\label{sec:progress}



\subsection*{Summary}

\begin{table}[H]
\centering
\small
\begin{tabular}{@{}clcc@{}}
\toprule
\textbf{WP} & \textbf{Work package} & \textbf{Target} & \textbf{Complete} \\
\midrule
WP1 & Foundations: prescribed textbook          & 31 Aug & \pct{55} \\
WP2 & Advanced sources                          & 14 Sep & \pct{30} \\
WP3 & Literature survey and Zotero              & 21 Sep & \pct{35} \\
WP4 & Writing: formal, historical, methodological chapters & 24 Aug & \pct{90} \\
WP5 & Experiment A: \code{herdtools7} litmus tests  & 28 Sep & \pct{10} \\
WP6 & Experiment B: C++11 microbenchmarks + \code{perf} & 19 Oct & \pct{0} \\
WP7 & Experiment C: gem5 (optional)               & 02 Nov & \pct{0} \\
WP8 & Final assembly and appendices             & 11 Nov & \pct{0} \\
\midrule
    & \textbf{Overall}                          &        & \pct{28} \\
\bottomrule
\end{tabular}
\caption{Work package status as of 24 August 2026.}
\end{table}

\noindent\fbox{\parbox{0.96\textwidth}{\small
\textbf{Note to the author:} every percentage in this section and in the Gantt
chart must be corrected to reflect what has \emph{actually} been done before
submission. Figures that overstate progress will be contradicted by the next
draft and are not worth the marks they might gain.}}

\subsection*{Gantt Chart}

\clearpage
\begin{figure}[H]
\centering
\begin{ganttchart}[
    hgrid, vgrid,
    x unit=7mm,
    y unit title=7mm,
    y unit chart=5mm,
    title height=1,
    bar height=0.6,
    group label font=\bfseries\scriptsize,
    bar label font=\scriptsize,
    milestone label font=\scriptsize\itshape,
    title label font=\scriptsize\bfseries,
    group/.append style={fill=wpcol, draw=wpcol},
    bar/.append style={fill=donecol!70, draw=donecol},
    progress label text={\scriptsize\textbf{#1\%}},
    progress label anchor/.style={left=2pt},
  ]{1}{13}

  \gantttitle{Aug}{2}
  \gantttitle{Sep}{4}
  \gantttitle{Oct}{4}
  \gantttitle{Nov}{3} \\
  \gantttitlelist{1,...,13}{1} \\

  \ganttgroup[progress=55]{WP1 Foundations}{1}{3} \\
  \ganttbar[progress=100]{T1.1 Sarangi Ch.\ 11}{1}{1} \\
  \ganttbar[progress=70]{T1.2 Sarangi \S12.1--12.4.6}{1}{2} \\
  \ganttbar[progress=20]{T1.3 Sarangi \S12.4.7}{2}{3} \\
  \ganttbar[progress=40]{T1.4 Sarangi \S10.11.4}{2}{3} \\

  \ganttgroup[progress=30]{WP2 Advanced sources}{1}{5} \\
  \ganttbar[progress=45]{T2.1 Next-Gen \S9.1--9.4}{1}{3} \\
  \ganttbar[progress=15]{T2.2 Next-Gen \S9.5--9.9}{3}{5} \\
  \ganttbar[progress=40]{T2.3 Primer Ch.\ 3--5}{1}{3} \\
  \ganttbar[progress=10]{T2.4 Primer Ch.\ 6--8, 11}{3}{5} \\

  \ganttgroup[progress=35]{WP3 Literature \& Zotero}{1}{6} \\
  \ganttbar[progress=80]{T3.1 Zotero library}{1}{2} \\
  \ganttbar[progress=30]{T3.2 Six core papers}{2}{4} \\
  \ganttbar[progress=5]{T3.3 Vendor manuals}{4}{6} \\

  \ganttgroup[progress=90]{WP4 Writing (Ch.\ 1--7)}{1}{2} \\
  \ganttbar[progress=100]{T4.1 Ch.\ 1--3}{1}{1} \\
  \ganttbar[progress=95]{T4.2 Ch.\ 4 formal}{1}{2} \\
  \ganttbar[progress=90]{T4.3 Ch.\ 5 history}{1}{2} \\
  \ganttbar[progress=85]{T4.4 Ch.\ 6 methodology}{1}{2} \\

  \ganttgroup[progress=10]{WP5 Experiment A (litmus)}{2}{7} \\
  \ganttbar[progress=25]{T5.1 herdtools7 setup}{2}{3} \\
  \ganttbar[progress=5]{T5.2 Litmus corpus}{3}{5} \\
  \ganttbar[progress=0]{T5.3 Run on x86; herd7 cross-check}{5}{6} \\
  \ganttbar[progress=0]{T5.4 Write Ch.\ 8}{6}{7} \\

  \ganttgroup[progress=0]{WP6 Experiment B (C++11)}{4}{11} \\
  \ganttbar[progress=0]{T6.1 Workload implementation}{4}{6} \\
  \ganttbar[progress=0]{T6.2 memory\_order variants}{5}{7} \\
  \ganttbar[progress=0]{T6.3 Emitted-assembly audit}{7}{9} \\
  \ganttbar[progress=0]{T6.4 perf measurement}{8}{10} \\
  \ganttbar[progress=0]{T6.5 Write Ch.\ 9--10}{9}{11} \\

  \ganttgroup[progress=0]{WP7 Experiment C (gem5, opt.)}{9}{12} \\
  \ganttbar[progress=0]{T7.1 gem5 setup}{9}{11} \\
  \ganttbar[progress=0]{T7.2 Parameter sweeps}{11}{12} \\

  \ganttgroup[progress=0]{WP8 Final assembly}{10}{13} \\
  \ganttbar[progress=0]{T8.1 Ch.\ 12--13}{10}{12} \\
  \ganttbar[progress=0]{T8.2 Appendices A--G}{11}{13} \\
  \ganttbar[progress=0]{T8.3 Revision passes}{12}{13} \\

  \ganttmilestone{v0.1 (24 Aug)}{2} \\
  \ganttmilestone{v0.2}{7} \\
  \ganttmilestone{v0.3}{11} \\
  \ganttmilestone{v1.0 (11 Nov)}{13}
\end{ganttchart}
\caption{Project schedule. Week 1 commences 17 August 2026; week 13 commences
9 November 2026. Shaded portions of each bar indicate completion. Milestone
dates for v0.2 and v0.3 are estimates pending confirmation.}
\end{figure}
\clearpage

\subsection*{Work Package Detail}

\paragraph{WP1 --- Foundations (\pct{55}).}
\begin{itemize}[leftmargin=*,itemsep=1pt]
\item \textbf{T1.1} Sarangi \cite{sarangi2025basic} Ch.~11, \emph{The Memory
  System}, \S11.1--11.3 (pp.~505--560). \pct{100}. Cache organisation, locality,
  and the latency figures that motivate the store buffer.
\item \textbf{T1.2} Sarangi Ch.~12, \S12.1--12.4.6 (pp.~565--603). \pct{70}.
  \S12.4.2 and \S12.4.3 read closely; \S12.4.6 in progress.
\item \textbf{T1.3} Sarangi \S12.4.7, \emph{Implementing a Memory Consistency
  Model} (pp.~603--608). \pct{20}. Begun; requires \S12.4.6 first.
\item \textbf{T1.4} Sarangi \S10.11.4, \emph{Out-of-Order Pipelines}
  (pp.~492--505). \pct{40}. Needed for the squash mechanism central to Q3.
\end{itemize}

\paragraph{WP2 --- Advanced sources (\pct{30}).}
\begin{itemize}[leftmargin=*,itemsep=1pt]
\item \textbf{T2.1} Sarangi \cite{sarangi2023nextgen} \S9.1--9.4. \pct{45}.
  \S9.4 supplies the execution-witness and access-graph notation adopted
  throughout the paper, chosen over the alternatives for continuity with the
  course.
\item \textbf{T2.2} Sarangi \S9.5--9.9 (coherence protocols, memory models,
  race detection, transactional memory). \pct{15}.
\item \textbf{T2.3} Nagarajan et al.\ \cite{nagarajan2020primer} Ch.~3--5
  (SC, TSO, relaxed models). \pct{40}.
\item \textbf{T2.4} Nagarajan et al.\ Ch.~6--8, 11 (coherence protocols,
  specification and validation). \pct{10}.
\end{itemize}

\paragraph{WP3 --- Literature and Zotero (\pct{35}).}
\begin{itemize}[leftmargin=*,itemsep=1pt]
\item \textbf{T3.1} Zotero library created; 22 entries filed with PDFs.
  \pct{80}. Exported to \texttt{refs.bib}, under version control.
\item \textbf{T3.2} Six papers identified as load-bearing for the argument and
  to be read in full rather than consulted: Lamport
  \cite{lamport1979multiprocessor}; Dubois, Scheurich and Briggs
  \cite{dubois1986memory}; Adve and Hill \cite{adve1990weak}; Adve and
  Gharachorloo \cite{adve1996shared}; Gniady, Falsafi and Vijaykumar
  \cite{gniady1999sc}; Alglave, Maranget and Tautschnig
  \cite{alglave2014herding}. \pct{30}.
\item \textbf{T3.3} Vendor manuals \cite{intel_sdm,arm_arm,riscv_spec}.
  \pct{5}.
\end{itemize}

\noindent Further references filed and cited: Gharachorloo et al.\
\cite{gharachorloo1990memory}; Hill \cite{hill1998multiprocessors}; Manson,
Pugh and Adve \cite{manson2005java}; Boehm and Adve
\cite{boehm2008foundations}; Ceze et al.\ \cite{ceze2007bulksc}; Blundell,
Martin and Wenisch \cite{blundell2009invisifence}; Marino et al.\
\cite{marino2011sc}. Supporting texts: Hennessy and Patterson
\cite{hennessy2019quantitative}; Culler, Singh and Gupta
\cite{culler1998parallel}; Herlihy et al.\ \cite{herlihy2020art}.

\paragraph{WP4 --- Writing (\pct{90}).}
A 42-page draft of the paper proper exists in the repository
(\texttt{main.tex}), comprising:

\begin{table}[H]
\centering
\small
\begin{tabular}{@{}rlc@{}}
\toprule
\textbf{Ch.} & \textbf{Title} & \textbf{State} \\
\midrule
1 & Introduction --- motivating example, problem statement, contributions & drafted \\
2 & Background --- memory wall, caches, store buffer, out-of-order & drafted \\
3 & Cache Coherence, and What It Does Not Solve & drafted \\
4 & The Design Space of Memory Consistency Models & drafted \\
5 & Historical Evolution & drafted \\
6 & Methodology & drafted \\
7 & Source Texts and Their Coverage & drafted \\
\bottomrule
\end{tabular}
\caption{State of the paper draft. Seven chapters, 42 pages, five TikZ figures,
compiling without errors or undefined references.}
\end{table}

\noindent All figures are drawn in TikZ per the course requirement. The
document is built with \texttt{latexmk}; \texttt{main.pdf} is deliberately
excluded from version control and attached instead to tagged releases.

\paragraph{WP5 --- Litmus experiments (\pct{10}).}
Design complete and documented in Chapter~6 of the paper draft; implementation
begun. The harness must avoid \texttt{pthread\_barrier}, whose internal fences
suppress the window under observation, in favour of a sense-reversing spin
barrier built from relaxed atomics. Shared variables are padded to distinct
cache lines to prevent false sharing from confounding the coherence behaviour.

\paragraph{WP6--WP8 (\pct{0}).}
Not started; scheduled per the Gantt chart. Designs for WP6 (simulator) and
WP7 (verification) are specified in Chapter~6 of the paper draft, stated in
advance so that the methodology can be assessed independently of the results.

\subsection{Risks and Mitigations}

\begin{table}[H]
\centering
\small
\begin{tabular}{@{}p{4.6cm}cp{7.0cm}@{}}
\toprule
\textbf{Risk} & \textbf{Sev.} & \textbf{Mitigation} \\
\midrule
Litmus tests yield zero observations, making Ch.~8 empty
 & Med & Vary thread placement, optimisation level and inserted delay; a null
        result with the parameter sweep documented is itself reportable. \\
Simulator scope creep consumes the schedule
 & High & Feature list frozen after v0.2. Purpose-built simulator is the
         baseline; Tejas \cite{sarangi2023nextgen} evaluated in parallel as a
         stretch goal only, never a dependency. \\
No AArch64 hardware available for Q2
 & Med & Single-board computer, Android device under Termux, or institutional
         server; failing all three, Q2 narrows to x86 and the limitation is
         stated. \\
Topic requires material beyond the course
 & Low & Sarangi's companion volume \cite{sarangi2023nextgen} is pitched at
        exactly this level and is by the same author. \\
\bottomrule
\end{tabular}
\end{table}

\subsection{Research Questions and Experimental Plan}

The remainder of the paper turns from exposition to investigation.
The plan is stated now so that the methodology can be assessed in advance of
the results.

\subsubsection*{Research questions}

\begin{description}[leftmargin=1.6cm,style=nextline,font=\normalfont\bfseries]
\item[RQ1 Observation] Which ``surprising'' memory orderings are actually
  observable on commodity x86 hardware, and do they match the x86-TSO model
  exactly? (SB observable; IRIW not; MP safe without a fence.)
\item[RQ2 Cost] What is the concrete performance cost of each C++11
  \code{memory\_order} level on the same workload, and how does it map to the
  fences the compiler emits?
\item[RQ3 Microarchitecture] How do microarchitectural parameters ---
  store-buffer depth, speculation, coherence protocol --- change which
  orderings are observable and how expensive ordering is?
\end{description}

\subsubsection*{Experiment A --- litmus testing on real hardware (primary)}

Using \code{herdtools7} --- \code{litmus7} to run tests on hardware,
\code{herd7} to simulate a model --- run the canonical suite (SB, MP, LB, IRIW,
plus Dekker-style tests) on an x86 machine, collecting outcome histograms over
billions of iterations.
Cross-check every observed outcome against the \code{herd7} prediction under the
published x86-TSO \code{.cat} model.
Expected result: SB observable, IRIW never observed, MP safe without barriers,
confirming TSO.
This answers RQ1.

Adopting \code{herdtools7} rather than a hand-written harness is a deliberate
methodological choice: it is the standard tool of the field, it already
implements the affinity and stress options needed to observe rare outcomes, and
\code{herd7} provides the model-checking side without our having to build one.

\subsubsection*{Experiment B --- C++11 atomics microbenchmarks (primary)}

Implement small representative workloads --- a shared counter, a
producer/consumer following the MP pattern, and a spinlock handoff --- templated
over \code{memory\_order} (relaxed, acquire/release, seq\_cst).
Compile at a fixed optimisation level, inspect the emitted assembly to record
which fences each variant produces, and measure throughput and latency with
Linux \code{perf}.
This answers RQ2 and ties source semantics to hardware behaviour to measured
cost.

\subsubsection*{Experiment C --- gem5 simulation (optional)}

If time allows, use gem5 \cite{binkert2011gem5} to vary store-buffer depth and
coherence settings and observe the effect on litmus outcomes and microbenchmark
cost, isolating causes that are fixed on real silicon.
This addresses RQ3.
It is flagged optional given its setup complexity; the paper stands on
Experiments~A and~B.

\subsection{Immediate Plan to v0.2}

\begin{enumerate}[leftmargin=*,itemsep=2pt]
\item Complete WP1 and advance WP2; read the six load-bearing papers in full.
\item Install and validate \code{herdtools7}; reproduce one known litmus result.
\item Run the canonical suite on x86 and record outcome histograms.
\item Implement the three Experiment~B workloads and capture emitted assembly.
\item Draft the Experiment~A results chapter.
\item Merge the second author's survey draft into the paper proper per
  \code{MERGE-PLAN.md}.
\end{enumerate}

\subsection{Repository}

All work is under version control at\\
\texttt{https://github.com/y4shv/ELL305-term-paper} (private; access granted to
the instructor on request).

\medskip
\noindent Draft \#1 corresponds to tag \texttt{v0.1}.
\texttt{draft1.tex} builds this progress report; \texttt{main.tex} builds the
paper draft described in WP4.
\texttt{refs.bib} is exported from Zotero.
